% !TEX program = xelatex
\documentclass[12pt]{ctexart}

\usepackage[a4paper,margin=1in]{geometry}
\usepackage{amsmath,amssymb,mathtools,bm}
\usepackage{booktabs,longtable,tabularx,array}
\usepackage{enumitem}
\usepackage{setspace}
\usepackage{hyperref}
\usepackage{xcolor}
\usepackage{titlesec}

\setstretch{1.22}
\setlength{\parindent}{2em}
\setlength{\parskip}{0.25em}
\setlist[itemize]{leftmargin=2em,itemsep=0.2em,topsep=0.3em}
\setlist[enumerate]{leftmargin=2em,itemsep=0.25em,topsep=0.35em}
\emergencystretch=2.5em
\sloppy

\newcolumntype{Y}{>{\raggedright\arraybackslash}X}
\newcolumntype{L}[1]{>{\raggedright\arraybackslash}p{#1}}

\titleformat{\section}{\normalfont\Large\bfseries}{\thesection.}{0.7em}{}
\titleformat{\subsection}{\normalfont\large\bfseries}{\thesubsection.}{0.6em}{}
\titleformat{\subsubsection}{\normalfont\normalsize\bfseries}{\thesubsubsection.}{0.5em}{}

\title{第17页说明：数据、改革前暴露度构造与类型结果\\\large 变量构造、经验设计、机制解释与文献依据}
\author{}
\date{2026年4月29日}

\begin{document}
\maketitle

\section*{核心结论}

第17页的核心不是说“改革前已经存在 MAH 制度意义上的 type-B 企业”。更准确的理解是：
\begin{center}
\fbox{\parbox{0.90\textwidth}{\centering 用改革前的企业或地区特征，构造其对 MAH 所解决的 implementation bottleneck 的暴露度。}}
\end{center}
也就是说，经验设计要识别的是：哪些企业或地区在改革前更受“研发能力与生产能力不能有效分离”这一瓶颈约束；如果 MAH 的主要机制是降低外部商业化的治理与合规摩擦，那么改革后，这些高暴露企业或地区应当出现更强的原研药实现、外部商业化路线使用、研发型进入或组织分化变化。

本文建议将该页标题从 \textit{Data, Exposure Construction, and Type Outcomes} 改为更清楚的：
\[
\boxed{\textit{Data, Pre-Reform Exposure, and Mechanism Outcomes}.}
\]
这能避免误解：type outcomes 是机制验证的一部分，不是论文的最终结果对象。论文最终对象仍然应是原研药创新与原研药实现。

\section{为什么这页重要}

当前模型的核心机制可以概括为：
\[
\tau_{ext}^{post}<\tau_{ext}^{pre},
\]
其中 $\tau_{ext}$ 表示研发者保留 holder 地位、通过外部合格生产商完成商业化时的治理与合规摩擦。MAH 不是需求冲击，不是科学生产率冲击，也不是所有企业共同收益的泛化平移；它改变的是原研药 idea 从 invention 到 commercialization 的组织路线。

因此，经验设计不能只写一个简单的改革前后平均变化：
\[
Y_{it}=\alpha_i+\gamma_t+\beta Post_t+\varepsilon_{it}.
\]
这种写法最多说明改革后总体是否变化，却不能说明变化是否来自 MAH 所解决的“外部商业化瓶颈”。更贴合机制的做法是：比较改革前更受 implementation bottleneck 约束的企业或地区，在改革后是否发生更强的变化。

这就是本页中 \textit{pre-reform exposure proxies} 的意义。

\section{什么是 exposure}

这里的 exposure 不是 treatment dummy，也不是改革后的 MAH-holder 身份。错误写法是：
\[
Exposure_i=1\{\text{企业改革后成为 MAH holder}\}.
\]
这个变量明显内生，因为改革后成为 holder 本身可能就是政策影响的结果。高能力企业、融资更强企业、增长更快企业更可能成为 holder；因此，用改革后 holder 身份定义 treatment 会把结果变量的一部分误当成处理变量。

正确的 exposure 应该是政策前确定的变量：
\[
Exposure_i^{pre}=\text{改革前企业对 MAH 所解决的商业化瓶颈的暴露程度}.
\]
MAH 主要解决的是：
\[
\text{有研发项目或原研 idea，但缺乏内部生产和商业化能力。}
\]
所以最核心的暴露度应写成：
\[
\boxed{
Exposure_i^{B,pre}=ResearchOrientation_i^{pre}\times WeakManufacturing_i^{pre}.
}
\]
其经济含义是：企业在改革前越像“有研发但缺生产”的主体，就越可能从 MAH 降低外部商业化摩擦中获益。

\section{三类 exposure 的构造}

第17页把暴露度分为三类：integrated-incumbent exposure、potential type-B exposure、type-C exposure。它们对应模型中的 $A/B/C$，但需要强调：这里的 $A/B/C$ 是经济学上的组织状态或经验分类，不是工商登记意义上的永久企业身份。

\subsection{Integrated-incumbent exposure：潜在 A 型暴露度}

页中写的是：
\begin{quote}
Integrated-incumbent exposure: original-drug activity plus relevant internal manufacturing capability.
\end{quote}
也就是：
\[
Exposure_i^{A,pre}=ResearchOrientation_i^{pre}\times ManufacturingCapability_i^{pre}.
\]
它表示改革前企业已经同时具备两类能力：第一，原研药或创新药相关活动；第二，内部生产能力。这类企业更接近 type-A integrated firm。

\subsubsection*{变量构造建议}

研发导向可以先用简单版本：
\[
ResearchOrientation_i^{pre}=zscore\left(\log(1+OriginalDrugApplications_{i,pre})\right).
\]
如果数据更丰富，可以构造综合指标：
\[
\begin{aligned}
ResearchOrientation_i^{pre}={}&w_1\log(1+IND_{i,pre}^{O})+w_2\log(1+Patent_{i,pre})\\
&+w_3R\&DIntensity_{i,pre}.
\end{aligned}
\]
制造能力可以定义为：
\[
\begin{aligned}
ManufacturingCapability_i^{pre}=zscore\big(&License_i^{pre}+GMP_i^{pre}+DosageScope_i^{pre}\\
&+InternalProducerRecords_i^{pre}\big).
\end{aligned}
\]
数据不足时，可以先用二元指标：
\[
ManufacturingCapability_i^{pre}=1\{\text{企业改革前有生产许可证或生产记录}\}.
\]
因此，潜在 A 型暴露度为：
\[
Exposure_i^{A,pre}=ResearchOrientation_i^{pre}\times ManufacturingCapability_i^{pre}.
\]

\subsubsection*{解释}

这类企业改革前已经能在企业内部把研发和生产连接起来，因此 MAH 对其直接瓶颈缓解效应可能弱于潜在 B 型企业。但它们仍然可能通过更灵活的生产安排、外包路线、管线扩展或行业均衡变化获益。

\subsection{Potential type-B exposure：潜在 B 型暴露度}

页中写的是：
\begin{quote}
Potential type-B exposure: research-oriented firms or entities with weak or no internal manufacturing capability.
\end{quote}
这是最重要的处理强度变量。推荐定义为：
\[
\boxed{
Exposure_i^{B,pre}=ResearchOrientation_i^{pre}\times WeakManufacturing_i^{pre}.
}
\]
其中：
\[
WeakManufacturing_i^{pre}=1-ManufacturingCapability_i^{pre},
\]
或者：
\[
WeakManufacturing_i^{pre}=1\{ManufacturingCapability_i^{pre}\leq q_{25}\}.
\]
如果要做更直观的经验分类，可以按研发导向和生产能力交叉分组。

\begin{table}[h!]
\centering
\caption{改革前研发导向与生产能力的经验分类}
\begin{tabularx}{\textwidth}{L{0.18\textwidth}L{0.18\textwidth}Y}
\toprule
研发导向 & 生产能力 & 经验解释 \\
\midrule
高 & 高 & 潜在 A 型：研发-生产一体化企业 \\
高 & 低 & 潜在 B 型：有研发导向但受内部生产约束的企业，核心处理组 \\
低 & 高 & 潜在 C 型：制造能力较强、研发导向较弱的生产专门化企业 \\
低 & 低 & 边缘企业或非核心样本 \\
\bottomrule
\end{tabularx}
\end{table}

\subsubsection*{为什么不能说“改革前 type-B 已经存在”}

MAH 意义上的 type-B 是制度承认后的组织路线：研发者可以保留上市许可，并委托合格生产商生产。改革前即使存在研发机构、科研团队、轻资产 biotech 或弱生产能力企业，也不等于它们已经拥有 MAH 制度下的 holder 路线。

因此，页中黄字非常关键：
\begin{quote}
Use pre-reform characteristics to measure exposure to the implementation bottleneck, not to claim that MAH-style type-B firms already existed before the reform.
\end{quote}
中文应写为：改革前的研发导向和弱生产能力只是用于衡量企业对“研发-生产绑定”瓶颈的暴露程度，而不是声称改革前已经存在 MAH 制度意义上的 type-B 企业。

\subsection{Type-C exposure：潜在 C 型暴露度}

页中写的是：
\begin{quote}
Type-C exposure: qualified manufacturing specialists that can support external implementation.
\end{quote}
这类企业不是主要的原研 idea 生成者，而是外部商业化路线中的生产供给方。可定义为：
\[
Exposure_i^{C,pre}=QualifiedManufacturing_i^{pre}\times (1-ResearchOrientation_i^{pre}).
\]
更简单的写法是：
\[
Exposure_i^{C,pre}=1\{\text{改革前有合格生产资质且研发活动较弱}\}.
\]

可使用的 proxy 包括：
\begin{itemize}
    \item 是否有药品生产许可证；
    \item 是否有 GMP 相关资质；
    \item 剂型覆盖范围；
    \item 历史生产批文数量；
    \item 是否在批文数据中频繁作为 NMPA 生产企业；
    \item 是否有委托生产或受托生产记录；
    \item 生产地址、剂型、产线范围。
\end{itemize}

\subsubsection*{Type-C 的经验预测}

Type-C 的主要 outcome 不应是“自己做更多原研药”，而应是：
\[
\text{external manufacturing demand and supply role increases.}
\]
改革后，高 type-C exposure 的企业可能更频繁出现在 NMPA 生产企业字段中，更多与 MAH holder 形成 holder-producer split，委托生产或受托生产记录增加，并在合格生产服务市场中承担更重要的生产端角色。

\section{核心实证设计}

\subsection{企业层面基准式}

最基本的 firm-level 设计是：
\[
Y_{it}=\alpha_i+\gamma_t+\beta\left(Post_t\times Exposure_i^{B,pre}\right)+X_i^{pre}\times\gamma_t+\varepsilon_{it}.
\]
其中：
\begin{itemize}
    \item $Y_{it}$：企业 $i$ 在年份 $t$ 的 outcome；
    \item $\alpha_i$：企业固定效应；
    \item $\gamma_t$：年份固定效应；
    \item $Post_t$：MAH 改革后；
    \item $Exposure_i^{B,pre}$：改革前潜在 B 型暴露度；
    \item $X_i^{pre}\times\gamma_t$：改革前企业规模、地区、治疗领域、所有制、历史创新能力等变量与年份交互，用于吸收差异化趋势。
\end{itemize}
核心系数 $\beta$ 的解释是：
\begin{center}
\fbox{\parbox{0.90\textwidth}{\centering 改革后，改革前更像“有研发但缺生产”的企业，是否出现更大的创新实现或外部商业化变化？}}
\end{center}

\subsection{更严格的 triple-difference 写法}

如果要直接体现“研发导向 $\times$ 弱生产能力”的互补关系，可以写：
\[
\begin{aligned}
Y_{it}={}&\alpha_i+\gamma_t\\
&+\beta\left(Post_t\times ResearchOrientation_i^{pre}\times WeakManufacturing_i^{pre}\right) \\
&+\theta_1\left(Post_t\times ResearchOrientation_i^{pre}\right)\\
&+\theta_2\left(Post_t\times WeakManufacturing_i^{pre}\right)+\varepsilon_{it}.
\end{aligned}
\]
这里 $\beta$ 不是简单说“研发型企业变多了”，也不是简单说“弱生产企业变多了”，而是说：
\begin{center}
\fbox{\parbox{0.88\textwidth}{\centering 改革后，既有研发导向、又缺内部生产能力的企业，变化是否更强？}}
\end{center}
这个设计最贴合机制。因为 MAH 不是单纯帮助所有研发企业，也不是单纯帮助所有没有厂的企业，而是帮助那些有原研项目但受内部生产约束限制的主体。

\subsection{地区层面写法}

如果企业层面数据不完整，可以先做 region-level exposure：
\[
Exposure_r^{B,pre}=\frac{\sum_{i\in r}ResearchOrientation_i^{pre}\times WeakManufacturing_i^{pre}}{\sum_{i\in r}1}.
\]
或者：
\[
Exposure_r^{B,pre}=\frac{\text{改革前高研发、弱生产企业数}}{\text{改革前医药企业总数}}.
\]
然后估计：
\[
Y_{rt}=\alpha_r+\gamma_t+\beta\left(Post_t\times Exposure_r^{B,pre}\right)+\varepsilon_{rt}.
\]
其中 $Y_{rt}$ 可以是地区层面的原研药批准数、原研药申请数、原研药批准占比、holder-producer split 比例、新进入研发型企业数、受托生产企业活跃度或 MAH holder 数量。

\section{Outcome 应该如何分层}

Type outcomes 不是最终唯一 outcome。当前论文主轴仍然是原研药创新。因此，建议把 outcomes 分成两层。

\subsection{第一层：最接近机制和最终对象的 outcomes}

第一层 outcome 应直接服务于“MAH 是否促进原研药实现”这一核心问题。建议包括：
\[
Y_{it}^{O}=\text{realized original-drug products},
\]
\[
Share_{it}^{O}=\frac{\text{original-drug realized products}}{\text{all realized products}},
\]
\[
RouteUse_{it}^{proxy,O}=\text{holder-producer split among original-drug records}.
\]
具体变量可以包括：
\begin{itemize}
    \item 原研药批准数；
    \item 原研药申请数；
    \item 原研药批准占比；
    \item CDE 企业名称与 NMPA 企业名称是否分离；
    \item MAH 企业与生产企业是否分离；
    \item 原研药记录中的 holder-producer split；
    \item 原研药从申请到批文的转化率或实现率。
\end{itemize}

\subsection{第二层：组织分化和 type outcomes}

第二层才对应 slide 中的 type outcomes：
\[
TypeA_{it},\quad TypeB_{it},\quad TypeC_{it}.
\]
可构造为：
\[
TypeA_{it}=1\{ResearchOrientation_{it}=1\ \&\ ManufacturingCapability_{it}=1\},
\]
\[
TypeB_{it}=1\{ResearchOrientation_{it}=1\ \&\ ManufacturingCapability_{it}=0\},
\]
\[
TypeC_{it}=1\{ResearchOrientation_{it}=0\ \&\ ManufacturingCapability_{it}=1\}.
\]
但必须谨慎表述：这些是 empirical classification，不是企业永久身份。一个企业可以在某一类产品上像 A，在另一类产品上像 B。更稳的写法是：
\begin{quote}
We classify firm-year or firm-product observations into A/B/C-like organizational states based on observable research and manufacturing characteristics.
\end{quote}
不要写成：
\begin{quote}
Firm $i$ is permanently a type-B firm.
\end{quote}

\section{为什么 post-reform R\&D effort 不能作为默认控制变量}

页中最后一句非常重要：
\begin{quote}
Post-reform R\&D effort is treated as a mechanism outcome, not as a default control variable.
\end{quote}
这是标准的 bad control 问题。理论链条是：
\[
\tau_{ext}\downarrow \Rightarrow H_B^{ext}\uparrow \Rightarrow \Psi_B\uparrow \Rightarrow x_B\uparrow \Rightarrow OriginalDrugRealization\uparrow.
\]
其中 $x_B$ 就是研发努力或原研方向投入。如果主回归控制改革后的 R\&D effort：
\[
Y_{it}=\alpha_i+\gamma_t+\beta Post_t\times Exposure_i^{pre}+\delta R\&D_{it}^{post}+\varepsilon_{it},
\]
就会把机制路径的一部分控制掉。这样估计出来的 $\beta$ 不是总效应，而是“在研发努力不变条件下的残余效应”。这会人为削弱政策通过提高原研方向研发投入而起作用的机制。

正确做法是把 post-reform R\&D effort 作为 outcome：
\[
R\&D_{it}=\alpha_i+\gamma_t+\beta_R(Post_t\times Exposure_i^{B,pre})+\varepsilon_{it},
\]
再看原研药结果：
\[
Y_{it}^{O}=\alpha_i+\gamma_t+\beta_Y(Post_t\times Exposure_i^{B,pre})+\varepsilon_{it}.
\]
如果要做机制分解，可以单独做 mediation-style analysis，但不应在主回归里默认控制 post-reform R\&D effort。

\section{为什么必须使用 pre-reform exposure}

识别上最关键的一点是：treatment intensity 必须在政策发生前确定。使用改革前变量：
\[
ResearchOrientation_i^{pre},\quad ManufacturingCapability_i^{pre},
\]
相当于问：在改革发生前，哪些企业最可能被“研发和生产必须绑定”这个制度约束卡住？

这种设计和 applied micro 中的 exposure design 或 heterogeneous treatment design 一致：用政策前暴露度刻画政策冲击的强弱，而不是用政策后行为定义处理组。它的优点是减少反向因果和选择内生性，避免把政策反应本身误当成外生处理。

\section{文献与理论依据}

\subsection{动态异质企业理论}

Hopenhayn (1992) 将异质企业、进入退出和长期稳态行业分布放在统一的动态行业均衡框架中。这个逻辑支持本文把改革前企业状态差异作为改革后创新、进入、退出和组织路线反应差异的基础。

Ericson and Pakes (1995) 的动态行业模型进一步强调企业投资、状态演化和退出决策共同塑造行业动态。本文虽然不必采用完整的 Ericson-Pakes 估计框架，但可以借用其状态变量、动态控制和企业价值函数语言。

\subsection{创新强度与异质创新边际}

Klette and Kortum (2004) 为 innovation intensity 与 firm dynamics 的结合提供了标准框架。Akcigit and Kerr (2018) 则强调不同类型创新和企业异质性对增长和行业创新结果的重要性。这支持本文区分 generic-oriented effort 与 original-drug-oriented effort，而不是只使用一个总研发变量。

\subsection{Invention 与 implementation 的分离}

Fons-Rosen, Roldan-Blanco and Schmitz (2024) 的核心启发在于：ideas 不是自动成为创新结果，而是需要 costly implementation。本文可借用这一分离逻辑，但不需要引入其并购、匹配或 Nash bargaining 模块。本文对应的 implementation 不是 acquisition，而是通过内部生产、委托生产、项目转让或放弃来完成原研药商业化。

\subsection{Complementary assets 与 commercialization}

Teece (1986) 说明，创新者能否获得创新收益取决于互补资产、商业化能力、整合或合作安排。这直接支持本文的经验暴露度构造：MAH 最应影响的是“有研发但缺互补生产/商业化资产”的主体，即：
\[
ResearchOrientation_i^{pre}\times WeakManufacturing_i^{pre}.
\]

\subsection{医药创新的 expected return 与制度设计文献}

Acemoglu and Linn (2004) 表明医药创新会对潜在市场规模和预期收益变化作出反应。Lakdawalla (2018) 系统讨论了药品行业中研发、定价、政策和激励之间的关系。Jia, Ma, Yang and Zhang (2023) 关于中国药品监管改革的研究表明，监管改进可以通过减少审批等待、提高预期回报、改变企业进入和创新结果来促进药品创新。本文与这一路径一致，但机制更具体：MAH 不是简单缩短审批等待，而是降低原研项目外部商业化路线的治理与合规摩擦。

\section{可直接写进论文的经验段落}

\subsection*{英文版}

\begin{quote}
We construct pre-reform exposure measures to capture firms' ex ante exposure to the commercialization bottleneck addressed by the MAH reform. The purpose of these measures is not to claim that MAH-style type-B firms already existed before the reform. Rather, the measures identify firms or regions for which the pre-reform linkage between authorization and internal manufacturing was more likely to constrain the realization of research-originated original-drug projects. Our baseline exposure measure combines pre-reform research orientation with weak internal manufacturing capability. This measure is predetermined with respect to the reform and captures the model-implied complementarity between original-drug idea generation and the lack of internal implementation capacity.
\end{quote}

\subsection*{中文版}

\begin{quote}
我们使用改革前企业特征构造 MAH 暴露度，用以衡量企业在改革前受到“研发与生产绑定”所形成的商业化瓶颈约束的程度。该变量并不意味着 MAH 制度意义上的 type-B 企业在改革前已经存在，而是用于识别那些在改革前已经具备一定研发导向、但缺乏内部生产和商业化能力的潜在受约束主体。基准暴露度定义为改革前研发导向与弱内部生产能力的交互项。该变量在政策发生前确定，因而避免了使用改革后 MAH-holder 身份所带来的选择内生性。
\end{quote}

\subsection*{基准式}

\[
Y_{it}=\alpha_i+\gamma_t+\beta(Post_t\times Exposure_i^{B,pre})+X_i^{pre}\times\gamma_t+\varepsilon_{it}.
\]

解释：系数 $\beta$ 衡量的是，改革前更暴露于 implementation bottleneck 的企业，在改革后是否出现更大的原研药实现、外部商业化路线使用或研发型进入变化。

\section{给老师汇报时的核心话术}

可以直接这样讲：

\begin{quote}
这一页不是在说改革前已经有 MAH 意义上的 B 类企业。我们只是用改革前的企业特征来衡量哪些企业更可能被“有研发但缺生产能力”这个商业化瓶颈卡住。MAH 如果真的降低的是外部商业化路线的治理与合规摩擦，那么改革后受益最大的应该不是所有企业，而是改革前有原研研发基础、但缺乏内部生产和商业化能力的企业。因此，我们用改革前研发导向和弱生产能力的交互项来构造 baseline exposure。这样做的好处是 treatment intensity 是政策前确定的，不是改革后的选择结果。之后实证上看这些高暴露企业或地区，在改革后是否出现更大的原研药实现、holder-producer split、外部路线使用、研发型进入和组织分化变化。post-reform R\&D effort 不能作为默认控制变量，因为它本身就是机制结果；如果控制掉，等于把政策通过提高原研方向研发投入的机制路径人为截断。
\end{quote}

\section*{参考文献}
\begin{enumerate}[leftmargin=2em]
    \item Acemoglu, Daron, and Joshua Linn. 2004. ``Market Size in Innovation: Theory and Evidence from the Pharmaceutical Industry.'' \textit{Quarterly Journal of Economics}, 119(3): 1049--1090.
    \item Akcigit, Ufuk, and William R. Kerr. 2018. ``Growth through Heterogeneous Innovations.'' \textit{Journal of Political Economy}, 126(4): 1374--1443.
    \item Ericson, Richard, and Ariel Pakes. 1995. ``Markov-Perfect Industry Dynamics: A Framework for Empirical Work.'' \textit{Review of Economic Studies}, 62(1): 53--82.
    \item Fons-Rosen, Christian, Pau Roldan-Blanco, and Tom Schmitz. 2024. ``The Effects of Startup Acquisitions on Innovation and Economic Growth.'' Working paper.
    \item Hopenhayn, Hugo A. 1992. ``Entry, Exit, and Firm Dynamics in Long Run Equilibrium.'' \textit{Econometrica}, 60(5): 1127--1150.
    \item Jia, Ruixue, Xiao Ma, Jianan Yang, and Yiran Zhang. 2023. ``Improving Regulation for Innovation: Evidence from China's Pharmaceutical Industry.'' NBER Working Paper No. 31976.
    \item Klette, Tor Jakob, and Samuel Kortum. 2004. ``Innovating Firms and Aggregate Innovation.'' \textit{Journal of Political Economy}, 112(5): 986--1018.
    \item Lakdawalla, Darius N. 2018. ``Economics of the Pharmaceutical Industry.'' \textit{Journal of Economic Literature}, 56(2): 397--449.
    \item Teece, David J. 1986. ``Profiting from Technological Innovation: Implications for Integration, Collaboration, Licensing and Public Policy.'' \textit{Research Policy}, 15(6): 285--305.
\end{enumerate}

\end{document}
